%!TeX root = ../libro.tex

% The following lines set the fonts used by the document.
% "fontspec" has been imported (by "ctexbook").

\IfFontExistsTF{XITS-Regular.otf}{%
    \IfFontExistsTF{XITS-Italic.otf}{%
        \IfFontExistsTF{FiraSans-SemiBold.otf}{%
            \IfFontExistsTF{FiraSans-SemiBoldItalic.otf}{%
                \setmainfont[
                    Path,
                    UprightFont = *-Regular.otf,
                    BoldFont = FiraSans-SemiBold.otf,
                    ItalicFont = *-Italic.otf,
                    BoldItalicFont = FiraSans-SemiBoldItalic.otf
                ]{XITS}
            }{}
        }{}
    }{}
}{}

\IfFontExistsTF{FiraSans-Light.otf}{%
    \IfFontExistsTF{FiraSans-SemiBold.otf}{%
        \IfFontExistsTF{FiraSans-LightItalic.otf}{%
            \IfFontExistsTF{FiraSans-SemiBoldItalic.otf}{%
                \setsansfont[
                    Path,
                    UprightFont = *-Light.otf,
                    BoldFont = *-SemiBold.otf,
                    ItalicFont = *-LightItalic.otf,
                    BoldItalicFont = *-SemiBoldItalic.otf
                ]{FiraSans}
            }{}
        }{}
    }{}
}{}

\IfFontExistsTF{Sarasa Mono SC Light}{%
    \IfFontExistsTF{Sarasa Mono SC Light Italic}{%
        \IfFontExistsTF{Sarasa Mono SC SemiBold}{%
            \IfFontExistsTF{Sarasa Mono SC SemiBold Italic}{%
                \setmonofont[
                    UprightFont = * Light,
                    BoldFont = * SemiBold,
                    ItalicFont = * Light Italic,
                    BoldItalicFont = * SemiBold Italic
                ]{Sarasa Mono SC}
                \PassOptionsToClass{fontset=none}{ctexbook}
                \setCJKsansfont[
                    UprightFont = * Light,
                    BoldFont = * SemiBold,
                    ItalicFont = * Light Italic,
                    BoldItalicFont = * SemiBold Italic
                ]{Sarasa Mono SC}
                \setCJKmonofont[
                    UprightFont = * Light,
                    BoldFont = * SemiBold,
                    ItalicFont = * Light Italic,
                    BoldItalicFont = * SemiBold Italic
                ]{Sarasa Mono SC}
                \IfFontExistsTF{Source Han Serif CN Light}{%
                    \setCJKmainfont[
                        UprightFont = Source Han Serif CN Light,
                        BoldFont = * SemiBold,
                        ItalicFont = * Light Italic,
                        BoldItalicFont = * SemiBold Italic
                    ]{Sarasa Mono SC}
                }{%
                    \setCJKmainfont[
                        UprightFont = * Light,
                        BoldFont = * SemiBold,
                        ItalicFont = * Light Italic,
                        BoldItalicFont = * SemiBold Italic
                    ]{Sarasa Mono SC}
                }
            }{}
        }{}
    }{}
}{
    \IfFontExistsTF{SarasaMonoSC-Light.ttf}{%
        \IfFontExistsTF{SarasaMonoSC-LightItalic.ttf}{%
            \IfFontExistsTF{SarasaMonoSC-SemiBold.ttf}{%
                \IfFontExistsTF{SarasaMonoSC-SemiBoldItalic.ttf}{%
                    \setmonofont[
                        Path,
                        UprightFont = *-Light.ttf,
                        BoldFont = *-SemiBold.ttf,
                        ItalicFont = *-LightItalic.ttf,
                        BoldItalicFont = *-SemiBoldItalic.ttf
                    ]{SarasaMonoSC}
                    \PassOptionsToClass{fontset=none}{ctexbook}
                    \setCJKsansfont[
                        Path,
                        UprightFont = *-Light.ttf,
                        BoldFont = *-SemiBold.ttf,
                        ItalicFont = *-LightItalic.ttf,
                        BoldItalicFont = *-SemiBoldItalic.ttf
                    ]{SarasaMonoSC}
                    \setCJKmonofont[
                        Path,
                        UprightFont = *-Light.ttf,
                        BoldFont = *-SemiBold.ttf,
                        ItalicFont = *-LightItalic.ttf,
                        BoldItalicFont = *-SemiBoldItalic.ttf
                    ]{SarasaMonoSC}
                    \IfFontExistsTF{SourceHanSerifCN-Light.otf}{%
                        \setCJKmainfont[
                            Path,
                            UprightFont = SourceHanSerifCN-Light.otf,
                            BoldFont = *-SemiBold.ttf,
                            ItalicFont = *-LightItalic.ttf,
                            BoldItalicFont = *-SemiBoldItalic.ttf
                        ]{SarasaMonoSC}
                    }{%
                        \setCJKmainfont[
                            Path,
                            UprightFont = *-Light.ttf,
                            BoldFont = *-SemiBold.ttf,
                            ItalicFont = *-LightItalic.ttf,
                            BoldItalicFont = *-SemiBoldItalic.ttf
                        ]{SarasaMonoSC}
                    }
                }{}
            }{}
        }{}
    }{}
}

% I redefine the emphasis style.
\let\emph\relax
\DeclareTextFontCommand{\emph}{\bfseries}

% These are packages that deal with mathematical formulae.
\usepackage{amsmath}
\usepackage{mathtools}
\usepackage[
    math-style=ISO,
    bold-style=ISO,
    mathrm=sym,
    mathbf=sym,
    partial=upright,
    warnings-off={mathtools-colon,mathtools-overbracket}
]{unicode-math}
\allowdisplaybreaks[3]

% I specify the mathematical font(s) used by the document.
\IfFontExistsTF{XITSMath-Regular.otf}{%
    \setmathfont[]{XITSMath-Regular.otf}
}{}

% % I specify the mathematical font(s) used by the document.
% \IfFontExistsTF{XITSMath-Regular.otf}{%
%     \setmathfont[]{XITSMath-Regular.otf}
% }{}
% \IfFontExistsTF{texgyretermes-math.otf}{%
% \setmathfont[range=bb/{latin,Latin}]{texgyretermes-math.otf}
% }{}
% \IfFontExistsTF{STIXTwoMath-Regular.otf}{%
%     \setmathfont[range={"2218-"2218}]{STIXTwoMath-Regular.otf}
% }{}

% These are packages that deal with theorem environments.
\usepackage{amsthm}
\usepackage{thmtools}
\usepackage{thm-restate}

\declaretheoremstyle[%
    spaceabove=2ex,%
    spacebelow=2ex,%
    headfont=\normalfont\bfseries,%
    notefont=\normalfont,%
    notebraces={(}{)},%
    bodyfont=\normalfont,%
    headindent=\parindent,%
    headpunct={},%
    postheadspace=1em%
]{TheoremStyle}
\declaretheoremstyle[%
    spaceabove=2ex,%
    spacebelow=2ex,%
    headfont=\normalfont\bfseries,%
    notefont=\normalfont,%
    notebraces={(}{)},%
    bodyfont=\normalfont,%
    headindent=\parindent,%
    headpunct={},%
    postheadspace=1em,%
    qed=\qedsymbol%
]{ProofStyle}

\numberwithin{equation}{chapter}
\declaretheorem[%
    style=TheoremStyle,%
    name=定理,%
    sibling=equation%
]{theorem}
\declaretheorem[%
    style=TheoremStyle,%
    name=原理,%
    sibling=theorem%
]{principle}
\declaretheorem[%
    style=TheoremStyle,%
    name=假定,%
    sibling=theorem%
]{assumption}

\declaretheorem[%
    style=TheoremStyle,%
    name=定义,%
    sibling=theorem%
]{definition}
\declaretheorem[%
    style=TheoremStyle,%
    name=例,%
    sibling=theorem%
]{example}
\declaretheorem[%
    style=TheoremStyle,%
    name=注,%
    numbered=no%
]{remark*}

% I redefine the proof environment,
% so that I can use thmtools to typeset a proof,
% but still have an optional argument,
% the title of the proof.
% \let\proof\relax
% \let\endproof\relax
\renewcommand{\proofname}{证}
\newcommand{\nomodepruvo}{证}
\renewcommand*{\qedsymbol}{证完.}
\declaretheorem[%
    style=ProofStyle,%
    name=\protect\nomodepruvo,
    numbered=no%
]{pruvo}
\renewenvironment{proof}[1][\proofname]{\renewcommand{\nomodepruvo}{#1}\begin{pruvo}}{\end{pruvo}}

% I copied the magic from "memoir", a document class
% for typesetting poetry, fiction, non-fiction,
% and mathematical works.
% According to the basic user manual of "memoir",
% TeX tries very hard to keep text lines justified
% while keeping the interword spacing as constant as possible,
% but sometimes fails and complains about an overfull hbox.
% The default mode for LaTeX typesetting is \fussy
% where the (variation of) interword
% spacing in justified text is kept to a minimum.
% Following the \sloppy declaration
% there may be a much looser setting of justified text.
% Additionally the class ("memoir")
% provides the \midsloppy declaration
% which allows a setting somewhere between \fussy and \sloppy.
% Using \midsloppy one will get
% fewer overfull lines compared with \fussy
% and fewer obvious large interword spaces than with \sloppy.
\makeatletter
\providecommand*{\midsloppy}{%
    \tolerance 5000%
    \hbadness 4000%
    \emergencystretch 1.5em%
    \hfuzz .1\p@%
    \vfuzz\hfuzz}
\makeatother

\makeatletter
\providecommand*{\leqnomode}{\tagsleft@true\let\veqno\@@leqno}
\providecommand*{\reqnomode}{\tagsleft@false\let\veqno\@@eqno}
\makeatother

% Some xeCJK settings.
% However, they only work in XeLaTeX.
% https://tex.stackexchange.com/questions/13172/detect-which-tex-engine-is-used
% This packages allows one to check if XeLaTeX is used
% and to do something if it is the case.
\usepackage{ifxetex}
\ifxetex
    \usepackage{xpatch}
    \ExplSyntaxOn
    % https://github.com/CTeX-org/ctex-kit/issues/636
    % The following lines apply the patch
    % (described in the link above)
    % only if the package date is after 2022-08-05
    % _and_ before 2022-08-31.
    % The bug was introduced in v3.9.1 (on 2022-08-05)
    % and fixed in dev (on 2022-08-31).
    % More information is available at:
    % https://github.com/CTeX-org/ctex-kit/commit/a841dabaac7c2dccacd1858ba96f6eb38adfb335
    % https://github.com/CTeX-org/ctex-kit/commit/aa5f00db05f12034845f8e5c4119971b5aae33ff
    \IfPackageAtLeastTF{xeCJK}{2022-08-05}{
        \IfPackageAtLeastTF{xeCJK}{2022-08-31}{}{
            \xpatchcmd \__xeCJK_check_for_xglue:
            {\__xeCJK_if_last_glue:TF}
            {\__xeCJK_if_last_glue:T}
            {}{\PatchFailed}
        }
    }{}
    \ExplSyntaxOff
    \xeCJKsetup{%
        % https://github.com/CTeX-org/ctex-kit/issues/636
        xCJKecglue = true,%
        CJKspace = true,%
        CheckSingle = true,%
    }
\fi

% This packages allows one to check if LuaLaTeX is used.
\usepackage{ifluatex}

% These are modifications that are
% to be inserted after \begin{document}.
\AtBeginDocument{%
    % This line switches off extra space after punctuation.
    \frenchspacing

    % This line kills as many overfull lines as possible.
    % The definition is elsewhere in this file.
    \midsloppy

    % This line changes the font of URLs.
    \renewcommand*{\UrlFont}{\ttfamily}

    % I change some mathematical symbols
    % because of my personal taste.
    \renewcommand*{\leq}{\leqslant}
    \renewcommand*{\geq}{\geqslant}
    \renewcommand*{\mathellipsis}{\cdots}
}

% % This package inserts figures.
% \usepackage{graphicx}
% % This package inserts and configures captions.
% \usepackage{caption}
% % This package inserts sub-figures.
% \usepackage{subcaption}
% % This line sets the location of the images.
% \graphicspath{{bildoj/}}

% The following lines make the width of a quotation wider.
\makeatletter
\renewenvironment{quotation}
{\list{}{\listparindent 2\ccwd%
        \itemindent\listparindent
        % \rightmargin\leftmargin
        \leftmargin\rightmargin
        \parsep\z@ \@plus\p@}%
    \item\relax}
{\endlist}
\makeatother

% The following lines start a new page with each section.
% https://tex.stackexchange.com/questions/9497/start-new-page-with-each-section
\providecommand*{\phantomsection}{}
\usepackage{titlesec}
\newcommand*{\sectionbreak}{\clearpage\phantomsection}

\makeatletter
\if@twoside % if the document is two-sided
    % https://latex.org/forum/viewtopic.php?t=6965
    \raggedbottom

    % One uses "\usepackage[pass]{geometry}"
    % to see the information of page margins
    % in the log file
    % (One had better do it in a separate file,
    % because the current file is so big
    % that it will take unnecessarily longer time.)
    \setlength\oddsidemargin{31pt}
    \setlength\evensidemargin{31pt}
    \setlength\marginparwidth{38pt}
\else % if the document is one-sided
    % I redefine \cleardoublepage.
    \def\cleardoublepage{\clearpage\ifodd\c@page\else
            \hbox{}\newpage\if@twocolumn\hbox{}\newpage\fi\fi}
\fi
\makeatother

% I define a shortcut for the following text
% that is so often used.
\newcommand*{\maldevigalegajxo}{%
    \begin{remark*}
        本节是选学内容.
        不学本节不影响理解任何必学内容
        (也就是, 未被声明为选学内容的节).
    \end{remark*}%
}

% Here are commands for displaying section numbers:
% \sekcio{##} (section ##)
% and \malneprasekcio{##} (optional section ##).
\providecommand*{\sekcio}[1]{#1}
\providecommand*{\malneprasekcio}[1]{\textit{\textsf{#1}}}

% This command marks the use of Esperanto
% (so that one just searches "\esperanto"
% to see all Esperanto expressions used in the book).
\providecommand{\esperanto}[1]{#1}

% This command marks the use of English
% (but I do not mark every use of English).
\providecommand{\angla}[1]{#1}

% This command marks the use of Pinyin.
\providecommand{\pinjino}[1]{\textsf{(#1)}}

% This package deals with headers and footers.
\usepackage{fancyhdr}
\pagestyle{fancy}
\fancyhf{}
\setlength{\headheight}{15pt}
\fancyhead[LO]{\nouppercase{\rightmark}}
\fancyhead[RO]{\thepage}
\fancyhead[RE]{\nouppercase{\leftmark}}
\fancyhead[LE]{\thepage}
\renewcommand*{\headrulewidth}{0pt}
\renewcommand*{\footrulewidth}{0pt}

% If a chapter has no section,
% I mark the header on the odd pages as well.
% One has to put the redefinition after fancyhdr,
% if it is imported.
\renewcommand*{\chaptermark}[1]{%
    \markboth{\CTEXifname{\CTEXthechapter\hspace{\ccwd}}{%
        }#1}{\CTEXifname{\CTEXthechapter\hspace{\ccwd}}{%
        }#1}%
}

% This package allows one to check whether it is a draft
% and to do something if it is the case
% or something else if it is not the case.
\usepackage{ifdraft}

\ifdraft{% draft case
    % \pagestyle{plain}

    % https://tex.stackexchange.com/questions/448296/
    % The following lines do not greatly compress
    % of the PDF document,
    % so complication is probably faster.
    \ifxetex
        \special{dvipdfmx:config z 0}
        \special{dvipdfmx:config C 0x40}
    \fi
    \ifluatex
        \pdfvariable compresslevel 0
        \pdfvariable objcompresslevel 0
    \fi
}%
{% final case
}

% https://tex.stackexchange.com/a/688793
% The following lines prepend an asterisk to optional sections.
\newcommand\EbleEstasAsterisko{}% check definable
\let\EbleEstasAsterisko\relax
\renewcommand\thesection{\EbleEstasAsterisko\thechapter.\arabic{section}}
\newcommand\KunAsteriskoEnEnhavtabelo{%
    \addtocontents{toc}%
    {\def\EbleEstasAsterisko{\let\EbleEstasAsterisko\relax\string\llap{*}}}%
}
\newcommand\SenAsteriskoEnEnhavtabelo{%
    \addtocontents{toc}%
    {\def\EbleEstasAsterisko{\let\EbleEstasAsterisko\relax}}%
}

% The package allows URL breaks at any alphanumerical character.
\usepackage{xurl}

% In appendix A in the old versions,
% a pair of vertical lines were used to denote a determinant.
% On 2024-07-31, I decided completely to banish
% the piece of notation.
% In order better to indicate the changes in the source code,
% I defined an empty command,
% which technically does nothing.
% Almost every (yeah, there are exceptions) occurrence of
% "\begin{vmatrix} ... \end{vmatrix}"
% was replaced by
% "\tridet \det \begin{bmatrix} ... \end{bmatrix} \tridet".
\providecommand*{\tridet}{}
\renewcommand*{\tridet}{}

% This is a package that deals with hyperlinks.
\usepackage[%
    breaklinks,
    hidelinks,
    %hyperindex,
    unicode
]{hyperref}

% % This package includes xmp metadata in the PDF document.
% \usepackage{hyperxmp}

\title{行列式入门}
\author{\esperanto{Determinantulo}}
\date{2025-06-28}

% Make aliases for \@title, \@author, \@date.
% Use \let.
\makeatletter
\let\latitolo\@title
\let\laauxtoro\@author
\let\ladato\@date
\makeatother

% the year
\newcommand*{\lajaro}{2025}

% Inspired by
% https://www.zhihu.com/question/641238134/answer/3377102718
\newcommand{\inscribe}[2]{
    \ifdeflength{\namelen}{}{\newlength{\namelen}}
    \ifdeflength{\datelen}{}{\newlength{\datelen}}
    \ifdeflength{\maxlen}{}{\newlength{\maxlen}}
    \settowidth{\namelen}{#1}
    \settowidth{\datelen}{#2}
    \ifdimgreater{\namelen}{\datelen}{\setlength{\maxlen}{\namelen}}{\setlength{\maxlen}{\datelen}}
    % [l], [c], [r] stand for
    % left, centre, right,
    % respectively.
    \hfill\begin{minipage}[r]{\maxlen}
        \makebox[\maxlen][r]{#1} \\
        \makebox[\maxlen][r]{#2}
    \end{minipage}
}

% % This package facilitates maintenance of
% % this multi-file project.
% \usepackage{subfiles}

% Here are packages that are conditionally imported
% so that less time is needed to build subfiles.

\ifSubfilesClassLoaded{}{%
    % This package inserts the table of bibliography.
    \usepackage[%
        backend=biber,%
        style=gb7714-2015,%
        % gbalign=gb7714-2015,%
        defernumbers=true,%
    ]{biblatex}
    % I use a null heading for each bibliography.
    \defbibheading{bibliography}[]{}

    % Multiple .bib files can also be added.
    % bibliographic entries
    % Note that "cxapitroj/bibliografiaj eroj.bib" is used
    % instead of "bibliografiaj eroj.bib"
    % because this file "agordoj.tex"
    % is "\input"-ed into "libro.tex",
    % which is not in the same folder as
    % bibliografiaj eroj.bib.
    \addbibresource{\subfix{cxapitroj/bibliografiaj eroj.bib}}
}

\ifSubfilesClassLoaded{}{%
    % This is a package that includes many logotypes about TeX.
    \usepackage{hologo}
}

\ifSubfilesClassLoaded{}{%
    % This package adds hyperlinks to the table of contents.
    % One uses "numbered" to number the bookmarks of the PDF file.
    \usepackage[numbered]{bookmark}
}
% % associated setting
% \ifSubfilesClassLoaded{%
%     % all hypertext options are turned off
%     \hypersetup{draft}
% }{}

\ifSubfilesClassLoaded{}{%
    \hypersetup{%
        pdfauthor = {Determinantulo},
        pdfcreator = {Amo},
        pdfkeywords = {determinant, linear algebra, mathematics, matrix},
        pdflang = {zh-CN},
        pdfcreationdate = {D:20250628062831Z},
        pdfmoddate = {D:20250628062831Z},
        pdfproducer = {Amo},
        pdfsubject = {Enkonduko al determinantoj},
        pdftitle = {行列式入门},
    }
}

\ifSubfilesClassLoaded{}{%
    % I use this package to make figures.
    \usepackage{tikz}
    \usetikzlibrary{calc,matrix}
}

\ifSubfilesClassLoaded{}{%
    % This package makes long tables.
    \usepackage{longtable}
    % % This is a cell that allows line breaks.
    % \providecommand*{\cxelo}[2][c]{%
    %     \begin{tabular}[#1]{@{}l@{}}#2\end{tabular}}
}
